\documentclass[11pt]{article}
\usepackage[margin=2.4cm]{geometry}
\usepackage{amsmath,amssymb,amsthm,booktabs}
\usepackage[capitalize]{cleveref}
\newtheorem{theorem}{Theorem}
\newtheorem{lemma}{Lemma}
\newtheorem{proposition}{Proposition}
\newtheorem{corollary}{Corollary}
\newtheorem{conjecture}{Conjecture}
\newtheorem*{remark}{Remark}
\newcommand{\vt}{v_2}
\title{Two Clocks: the $\pm1$ Meters, Depth Supply,\\ and Unconditional Escape Bounds}
\author{John A. Janik\\ \small\texttt{john.janik@gmail.com}}
\date{June 2026}
\begin{document}
\maketitle

\begin{abstract}
We complete the run--budget clock of the companion note: the $b{=}1$ runs
of a Syracuse valuation word meter $2$-adic approach to the ghost $-1$
(depth decrement exactly $1$ per step), and --- new here --- the $b{=}2$
runs meter approach to the ghost $+1$, the \emph{trivial cycle}, with
depth decrement exactly $2$ per step and a parity handoff. Every letter
is a depth reading against one of the three rational points
$-1,\,+1,\,-\tfrac13$, and the word parses uniquely into alternating
meter readings (machine-verified). Ascent decomposes exactly into a
two-clock budget. Acyclicity makes every reading consume a congruence
point \emph{without replacement}, giving unconditional results: a
divergent orbit upcrosses the level $2^k$ at most $2^{k-3}$ times; an
acyclic orbit dwells in a band of ratio $\le2$ at most $O(\log Y)$
consecutive steps (the $D=\tfrac12$ strip carries exactly \emph{one}
word per phase); hence a divergent orbit spends at most $O(X\log X)$
steps of its entire lifetime below height $X$. Target 2 is then
isolated as a \emph{joint} calibration conjecture for two scalar
channels --- the $-1$ depth law and the $+1$ \emph{parity} law --- and we
show the first channel alone cannot suffice: a natural-depth orbit can
still ascend through pure parity bias, along the $(1,1,2)$-channel whose
limit is the negative rational cycle $-19/11$.
\end{abstract}

\section{The trichotomy and the two clocks}

Throughout $S(x)=(3x+1)/2^{b}$, $b=\vt(3x+1)$, acting on odd integers.

\begin{lemma}[Letter trichotomy]\label{lem:tri}
For odd $x$, exactly one of the following holds:
\begin{itemize}
\item[(i)] $x\equiv3\pmod4$, i.e.\ $\vt(x+1)\ge2$; then $b=1$;
\item[(ii)] $x\equiv1\pmod8$, i.e.\ $\vt(x-1)\ge3$; then $b=2$;
\item[(iii)] $x\equiv5\pmod8$, i.e.\ $\vt(x-1)=2$; then $b\ge3$.
\end{itemize}
Uniformly $b=\vt\bigl(x+\tfrac13\bigr)$ in $\mathbb Z_2$: every letter is
a $2$-adic depth reading against one of $-1$, $+1$, $-\tfrac13$.
\end{lemma}

\begin{proof}
Write $3x+1=3(x+1)-2=3(x-1)+4$. If $\vt(x+1)\ge2$ then
$\vt(3(x+1))\ge2$ and subtracting $2$ gives $\vt=1$: $b=1$. If
$x-1=2^hu$, $u$ odd, $h\ge3$: $3x+1=4(3\cdot2^{h-2}u+1)$ with the inner
factor odd, so $b=2$. If $h=2$: $3x+1=4(3u+1)$ with $3u+1$ even, so
$b\ge3$. Finally $\vt(3x+1)=\vt(3)+\vt(x+\tfrac13)=\vt(x+\tfrac13)$.
\end{proof}

\begin{lemma}[$-1$ clock; companion note]\label{lem:minus}
If $g=\vt(x+1)\ge2$ then $S(x)+1=3(x+1)/2$: the depth decreases by
exactly $1$ per step, the maximal $b{=}1$ run from $x$ has length
exactly $g-1$, and $g\le\log_2(x+1)$ for $x>0$.
\end{lemma}

\begin{lemma}[$+1$ clock]\label{lem:plus}
If $h=\vt(x-1)\ge3$ then $b=2$ and
\[
S(x)-1=\frac{3(x-1)}{4}:
\]
the depth against $+1$ decreases by exactly $2$ per step. The maximal
$b{=}2$ run from $x$ has length exactly $\lfloor(h-1)/2\rfloor$, with
terminal depth $1$ if $h$ is odd --- the run hands off directly to a
$b{=}1$ run --- and terminal depth $2$ if $h$ is even --- the next letter
is $b\ge3$. Moreover $h\le\log_2(x-1)$ for $x>1$, and $b{=}2$ forever
$\iff x=1$: \emph{the $+1$ meter measures $2$-adic proximity to the
trivial cycle}.
\end{lemma}

\begin{proof}
$S(x)=(3x+1)/4$ by \cref{lem:tri}, and $S(x)-1=(3x-3)/4=3(x-1)/4$, so
$\vt(S(x)-1)=h-2$. The run continues while the current depth is $\ge3$,
giving length $\lfloor(h-1)/2\rfloor$ and terminal depth
$h-2\lfloor(h-1)/2\rfloor\in\{1,2\}$ according to parity; terminal
depth $1$ means $x\equiv3\pmod4$ (case (i)), terminal depth $2$ is case
(iii). Positivity gives the cap; $x=1$ is the fixed point of
$x\mapsto(3x+1)/4$.
\end{proof}

\begin{proposition}[Alternating parse]\label{prop:parse}
Every Syracuse word is reproduced exactly by reading only the meters:
emit $g-1$ ones and jump $x\mapsto3^{\,g-1}(x+1)/2^{\,g-1}-1$; emit
$\lfloor(h-1)/2\rfloor$ twos and jump
$x\mapsto3^{\lfloor(h-1)/2\rfloor}(x-1)/4^{\lfloor(h-1)/2\rfloor}+1$;
at depth-$2$ points emit the single letter $b=\vt(x+\tfrac13)\ge3$.
(Machine-verified on $2{,}000$ random trajectories of length $300$;
clock laws verified exactly on $20{,}000$ samples.)
\end{proposition}

\begin{proposition}[Two-clock budget identity]\label{prop:budget}
Over any orbit window of $N$ letters, writing $g_i$ for the $1$-run
depths, $h_j$ for the $2$-run depths and $b_\ell\ge3$ for the big
letters,
\[
\log_2\frac{x_N}{x_0}
=(\log_23-1)\sum_i(g_i-1)
-(2-\log_23)\sum_j\Bigl\lfloor\frac{h_j-1}{2}\Bigr\rfloor
-\sum_\ell(b_\ell-\log_23)
+E,
\]
with $0\le E\le N/(3x_{\min}\ln2)$. Ascent is therefore financed by
exactly three channels: deep $-1$ readings ($+0.585$ per unit run
length), shallow-and-odd $+1$ readings ($-0.415$ per $2$-letter), and
avoidance of big letters ($-(b-1.585)$ each).
\end{proposition}

\section{Depth supply: readings are without replacement}

\begin{lemma}[Depth supply]\label{lem:supply}
Let the forward orbit be acyclic (e.g.\ divergent). For $2\le g\le k+1$,
the number of times the orbit \emph{ever} visits a point
$x\in[2^k,2^{k+1})$ with $\vt(x+1)\ge g$ is at most
$\max(2^{k-g},1)$, and likewise for $\vt(x-1)\ge g$. In particular each
$b{=}1$ step permanently consumes one point of the
$\equiv3\ (\mathrm{mod}\ 4)$ pool of its dyadic band, and the unique
deepest $-1$ point of band $k$ is the truncated ghost $2^{k+1}-1$.
\end{lemma}

\begin{proof}
An acyclic orbit visits each integer at most once; the points of depth
$\ge g$ in the band form the residue class $-1 \bmod 2^g$ (resp.\
$+1$), of cardinality $2^{k-g}$ for $g\le k$ and $1$ for $g=k+1$.
\end{proof}

\begin{theorem}[Upcrossing bound]\label{thm:upcross}
A divergent orbit makes at most $2^{k-3}$ upward crossings of the level
$2^k$, ever.
\end{theorem}

\begin{proof}
Only $b{=}1$ steps ascend, and they multiply by less than $2$; so the
step crossing $2^k$ starts at some $x\in[2^{k-1},2^k)$ with
$x\equiv3\pmod4$. By acyclicity these starting points are pairwise
distinct, and the pool has $2^{k-3}$ elements. (Using
$x\ge(2^{k+1}-1)/3$ refines the constant to $2^k/12+O(1)$.)
\end{proof}

\section{Logarithmic dwell and the escape rate}

\begin{lemma}[The dyadic strip carries one word]\label{lem:oneword}
For $D=\tfrac12$ the confinement window
$[j\log_23-D,\,j\log_23+D]$ has length $1$ and contains exactly one
integer for every $j$ (irrationality of $\log_23$); hence the exact
count of the parent note satisfies $N_{1/2}(p)=1$ for all $p$: the
unique confined word is the Sturmian word of slope $\log_23$.
(Verified by the exact big-integer recursion to $p=1600$.) Allowing an
arbitrary phase, the number of distinct length-$p$ blocks appearing in
any width-$1$-confined word is at most $p+1$, by the three-distance
argument.
\end{lemma}

\begin{theorem}[Logarithmic dwell at ratio $\le2$]\label{thm:logdwell}
An acyclic positive orbit remains within a band $[Y,WY]$, $W\le2$, for
at most $3\log_2Y+O(1)$ consecutive steps.
\end{theorem}

\begin{proof}
Set $p=\lceil\log_2 Y\rceil$. The segment's word is confined to a strip
of width $1+O(M/Y)$, so by \cref{lem:oneword} it has at most $2(p+1)$
distinct length-$p$ blocks (the $O(M/Y)$ widening can double the
integer choice at no more than one position once $Y\gg M$). A segment
of length $M>2(p+1)+p$ contains two positions carrying the same
length-$p$ block; both lie in the band, so they differ by less than
$Y\le2^{p}<2^{p+1}$, and repetition rigidity (parent program,
Lean-verified) forces the two points to be \emph{equal} --- an integer
cycle, contradicting acyclicity. Hence $M\le3p+O(1)$.
\end{proof}

\begin{theorem}[Escape rate]\label{thm:escape}
A divergent orbit spends at most $O(X\log X)$ steps of its entire
lifetime at heights $\le X$.
\end{theorem}

\begin{proof}
Visits to the dyadic band $[2^k,2^{k+1})$ are bounded by upcrossings of
$2^k$ plus downcrossings of $2^{k+1}$ plus one; downcrossings of a
level exceed its upcrossings by at most one, so by \cref{thm:upcross}
the band is visited at most $2^{k-3}+2^{k-2}+2\le2^{k-1}$ times. Each
visit lasts at most $3k+O(1)$ steps by \cref{thm:logdwell}. Summing
$\sum_{k\le\log_2X}2^{k-1}(3k+O(1))=O(X\log X)$.
\end{proof}

\begin{remark}
\cref{thm:escape} says divergence, if it occurs, is \emph{brisk}: the
orbit cannot oscillate indefinitely through bounded regions; its
occupation time below every height is finite and near-linear. All
ingredients are elementary given repetition rigidity.
\end{remark}

\section{Target 2 isolated: the joint calibration conjecture}

Under the natural ($2$-adic Haar) law the meters read
$\mathbb P[g=j]=2^{-(j-1)}$ ($j\ge2$, mean $3$),
$\mathbb P[h=j]=2^{-(j-1)}$ ($j\ge2$, mean $3$,
$\mathbb P[h\text{ odd}]=\tfrac13$), and the budget of
\cref{prop:budget} gives the familiar drift $\log_23-2<0$: descent.
Measured on random trajectories (windows of $60$ letters; ascending
$=$ end $\ge4\times$ start, never below start):

\begin{center}
\begin{tabular}{lccc}
\toprule
channel & natural & typical windows & ascending windows\\
\midrule
$-1$ depth mean $\mathbb E[g]$ & $3.000$ & $3.053$ & $3.990$\\
$+1$ depth mean $\mathbb E[h]$ & $3.000$ & $3.020$ & $3.196$\\
$+1$ parity $\mathbb P[h\ \mathrm{odd}]$ & $0.333$ & $0.358$ & $\mathbf{0.630}$\\
big letters per step & $0.250$ & $0.235$ & $0.093$\\
\bottomrule
\end{tabular}
\end{center}

\noindent The exact budget decomposition of the ascending windows
($+23.82$ from $1$-runs, $-5.69$ from $2$-letters, $-11.17$ from big
letters $=+6.95$, matching the measured net exactly) shows ascent drawing
on \emph{both} clocks: deeper $-1$ readings \emph{and} an odd-parity
bias of the $+1$ readings, the latter being precisely avoidance of big
letters after $2$-runs.

\begin{conjecture}[Two-clock calibration; Target 2]\label{conj:cal}
Along any divergent positive orbit the empirical meter laws calibrate
to natural: the run-start depth law tends to $2^{-(g-1)}$ \emph{and}
the handoff parity tends to $\tfrac13$. (Either statement is to be
read in Ces\`aro form; by \cref{prop:budget} joint calibration forces
mean drift $\log_23-2<0$, contradicting divergence.)
\end{conjecture}

\begin{remark}[The $g$-law alone cannot suffice]
A counterexample cannot be excluded by the $-1$ channel only: an orbit
with perfectly natural depth law ($\mathbb E[g]=3$) but pure parity
bias ($h\equiv3$ always) follows the letter pattern $(1,1,2)^\infty$,
which has mean $\tfrac43<\log_23$ and ascends. This channel's limit
object is the \emph{negative rational cycle}
$x=19/(2^4-3^3)=-19/11$, a ghost: the conjecture's parity half is
exactly the statement that positive integers cannot ride the
$-19/11$-ghost forever. Any proof of Target 2 must therefore meter
both clocks.
\end{remark}

\begin{remark}[Bridge to the mixed $(2,3)$-adic tower]
After a $1$-run of length $\ell=g-1$ from $x$ with $x+1=2^g m$ ($m$
odd), the next handoff reading is
\[
h \;=\; 1+\vt\bigl(3^{\ell}m-1\bigr):
\]
the $+1$ meter literally reads the $2$-adic valuation of
$3^{\ell}m-1$. Its parity --- the channel ascending orbits bias --- is a
joint condition on powers of $3$ against the \emph{higher} bits $m$ of
the current point, the lifting-the-exponent regime (for the truncated
ghosts $m=1$ the handoffs are fully deterministic:
$\vt(3^\ell-1)=1$ for $\ell$ odd, $=2+\vt(\ell)$ for $\ell$ even).
This places the band-crossing certificate of Target 2 squarely in the
joint $2^a3^b$ tower of the parent program: the certificate must meter
$\vt(3^\ell m-1)$, not $m$ alone.
\end{remark}

\begin{remark}[What remains, honestly]
The supply bounds of \cref{lem:supply} are permanent but the pools
regenerate with height: a divergent orbit gains fresh deep points as
it climbs, so counting alone cannot close \cref{conj:cal}. What the
two-clock structure changes is the \emph{shape} of the wall: quenched
normality of the full word is no longer needed --- only two scalar
renewal channels, each a congruence reading at handoff points, with
the $(2,3)$-tower coupling above as the natural attack surface.
\end{remark}

\end{document}
